
\begin{frame}{XPath}
XPath est un langage \alert{d'expression de chemins}.

\vfill

Il permet l'accès à des 
 n{\oe}uds d’un documents XML en décrivant les chemins menant ces n{\oe}uds.

\vfill

XPath est un outil de base dans la manipulation de documents XML. Utilisé dans 
\begin{description}
\item[XQuery] XPath fait partie du langage d’interrogation XQuery,
\item[XSLT] pour la sélection de la partie des données à transformer,
\item[XML Schema] pour les expressions de contraintes d’unicité,
%\item[XSL] pour la sélections des parties des données à formater,
\item[XPointer] pour l'identification de fragments,
\item[XLink] pour ancrer les liens hypertextes,
\item[...]
\end{description}
\end{frame}


\begin{frame}{XPath}
\frametitle{Les versions de XPath}

Toutes normalisées par le W3C: cf \url{http://www.w3c.org} pour les normes (bon courage!).

\begin{itemize}
\item XPath 1.0 est utilisé dans XSLT 1.0 (que nous \textit{pourrions} étudier dans ce cours),
\item XPath 1.0 est inclus dans XPath 2.0,
\item XPath 2.0 est inclus dans XQuery (que nous survolerons dans ce cours).
\end{itemize}

Nous commen\c{c}ons donc tout naturellement par étudier XPath 1.0.
\end{frame}


\subsection{Les chemins de localisation}

\begin{frame}{Une intuition}

Permet de ``pointer'' (localiser) des n{\oe}uds dans un document XML.

\begin{exampleblock}{Intuition}
Par \alert{navigation} : en se déplaçant dans l'arbre de "étape par étape" à partir de la racine, on va chercher à 
   atteindre certaines de ses parties.
\end{exampleblock}

\vfill

\begin{block}{Résultat d'une expression XPath}
Le résultat de l'évaluation d'une expression XPath est un \alert{ensemble de n{\oe}uds} sélectionnés dans le
document en entrée.
\end{block}

\vfill

\alert{Attention}.  Un n{\oe}ud peut être racine d'un sous-arbre; dans ce cas
c'est l'ensemble du sous-arbre qui fait partie du résultat.
 
\end{frame}

\begin{frame}[allowframebreaks]{Commen\c{c}ons par un exemple}

\begin{center}
\begin{minipage}{0.7\textwidth}
\lstinputlisting[basicstyle=\small, language=XML,breaklines=true, caption={Un simple fichier XML}]{codes/livresExXpathSansDTD.xml}
\end{minipage}
\end{center}

\framebreak

Expression XPath\,: \texttt{/child::library/child::book/child::author}

\begin{center}
\begin{tikzpicture}[auto]
\tiny
  \tikzstyle{emptyStyle}=[text width=0.9cm, text centered]
  \tikzstyle{parser}=[rounded corners=3mm,thick,draw=DarkViolet!75,fill=DarkViolet!20,,text centered]
  \tikzstyle{appli}=[rounded corners=1mm,thick,draw=red!75,fill=red!20,text centered,text width=2cm]
  \tikzstyle{every label}=[black]

  \begin{scope}
    \node[emptyStyle](r) at (6,5) {Root};
   \node[emptyStyle](comment) at (1,4) {\texttt{\verb?<!-- ceci ... >?}};
   \draw[-] (r) -- (comment);
    %\node[emptyStyle](entete) at (3,4) {\texttt{$<?xml$ version $...>$}};
   %\draw[-] (r) -- (entete);
    \node[emptyStyle](library) at (6,4) {\verb?<library> ... </library>?};
   \draw[-] (r) -- (library);

    \node[emptyStyle](book1) at (3,3) {\verb?<book> ... </book>?};
   \draw[-] (library) -- (book1);
    \node[emptyStyle](book2) at (6,3) {\verb?<book> ... </book>?};
   \draw[-] (library) -- (book2);
    \node[emptyStyle](book3) at (10,3) {\verb?<book> ... </book>?};
   \draw[-] (library) -- (book3);

    \node[emptyStyle](year1) at (2,3) {\verb?year = "1857"?};
   \draw[-] (book1) -- (year1);
    \node[emptyStyle](author1) at (2,2) {\verb?<author> ... </author>?};
   \draw[-] (book1) -- (author1);
    \node[emptyStyle](textauthor1) at (2,1) {\it Charles Baudelaire};
   \draw[-] (author1) -- (textauthor1);
    \node[emptyStyle](title1) at (3,2) {\verb?<title> ... </title>?};
   \draw[-] (book1) -- (title1);
    \node[emptyStyle](texttitle1) at (3,1) {\it Les fleurs du mal};
   \draw[-] (title1) -- (texttitle1);

    \node[emptyStyle](year2) at (5,3) {\verb?year = "1862"?};
   \draw[-] (book2) -- (year2);
    \node[emptyStyle](author2) at (6,2) {\verb?<author> ... </author>?};
   \draw[-] (book2) -- (author2);
    \node[emptyStyle](textauthor2) at (6,1) {\it Victor Hugo};
   \draw[-] (author2) -- (textauthor2);
    \node[emptyStyle](ln2) at (5,2) {\verb?ln = "Besancon"?};
   \draw[-] (author2) -- (ln2);
    \node[emptyStyle](title2) at (7,2) {\verb?<title> ... </title>?};
   \draw[-] (book2) -- (title2);
    \node[emptyStyle](texttitle2) at (7,1) {\it Les miserables};
   \draw[-] (title2) -- (texttitle2);

    \node[emptyStyle](year3) at (9,3) {\verb?year = "1877"?};
   \draw[-] (book3) -- (year3);
    \node[emptyStyle](author3) at (10,2) {\verb?<author> ... </author>?};
   \draw[-] (book3) -- (author3);
    \node[emptyStyle](textauthor3) at (10,1) {\it Emile Zola};
   \draw[-] (author3) -- (textauthor3);
    \node[emptyStyle](ln3) at (9,2) {\verb?ln = "Paris"?};
   \draw[-] (author3) -- (ln3);
    \node[emptyStyle](title3) at (11,2) {\verb?<title> ... </title>?};
   \draw[-] (book3) -- (title3);
    \node[emptyStyle](texttitle3) at (11,1) {\it L'assomoir};
   \draw[-] (title3) -- (texttitle3);

\end{scope}

%  \begin{pgfonlayer}{background}
%    \filldraw [line width=4mm,join=round,black!20]
%      (parser.north  -| parser.west)  rectangle (serial.south  -| serial.east);
%\end{pgfonlayer}


\end{tikzpicture}
\end{center}

\framebreak

On part de la racine\,: 
\texttt{\textcolor{red}{\bf /}child::library/child::book/child::author}

\begin{center}
\begin{tikzpicture}[auto]
\tiny
  \tikzstyle{emptyStyle}=[text width=0.9cm, text centered]
  \tikzstyle{red}=[rounded corners=1mm,thick,draw=red!75,fill=red!20,text width=0.9cm, text centered]
  \tikzstyle{DarkViolet}=[rounded corners=1mm,thick,draw=DarkViolet!75,fill=DarkViolet!20,minimum size=7mm,text width=0.8cm, text centered]
 \tikzstyle{every label}=[black]

  \begin{scope}
    \node[red](r) at (6,5) {Root};
   \node[emptyStyle](comment) at (1,4) {\texttt{\verb?<!-- ceci ... >?}};
   \draw[-] (r) -- (comment);
    %\node[emptyStyle](entete) at (3,4) {\texttt{$<?xml$ version $...>$}};
   %\draw[-] (r) -- (entete);
    \node[emptyStyle](library) at (6,4) {\verb?<library> ... </library>?};
   \draw[-] (r) -- (library);

    \node[emptyStyle](book1) at (3,3) {\verb?<book> ... </book>?};
   \draw[-] (library) -- (book1);
    \node[emptyStyle](book2) at (6,3) {\verb?<book> ... </book>?};
   \draw[-] (library) -- (book2);
    \node[emptyStyle](book3) at (10,3) {\verb?<book> ... </book>?};
   \draw[-] (library) -- (book3);

    \node[emptyStyle](year1) at (2,3) {\verb?year = "1857"?};
   \draw[-] (book1) -- (year1);
    \node[emptyStyle](author1) at (2,2) {\verb?<author> ... </author>?};
   \draw[-] (book1) -- (author1);
    \node[emptyStyle](textauthor1) at (2,1) {\it Charles Baudelaire};
   \draw[-] (author1) -- (textauthor1);
    \node[emptyStyle](title1) at (3,2) {\verb?<title> ... </title>?};
   \draw[-] (book1) -- (title1);
    \node[emptyStyle](texttitle1) at (3,1) {\it Les fleurs du mal};
   \draw[-] (title1) -- (texttitle1);

    \node[emptyStyle](year2) at (5,3) {\verb?year = "1862"?};
   \draw[-] (book2) -- (year2);
    \node[emptyStyle](author2) at (6,2) {\verb?<author> ... </author>?};
   \draw[-] (book2) -- (author2);
    \node[emptyStyle](textauthor2) at (6,1) {\it Victor Hugo};
   \draw[-] (author2) -- (textauthor2);
    \node[emptyStyle](ln2) at (5,2) {\verb?ln = "Besancon"?};
   \draw[-] (author2) -- (ln2);
    \node[emptyStyle](title2) at (7,2) {\verb?<title> ... </title>?};
   \draw[-] (book2) -- (title2);
    \node[emptyStyle](texttitle2) at (7,1) {\it Les miserables};
   \draw[-] (title2) -- (texttitle2);

    \node[emptyStyle](year3) at (9,3) {\verb?year = "1877"?};
   \draw[-] (book3) -- (year3);
    \node[emptyStyle](author3) at (10,2) {\verb?<author> ... </author>?};
   \draw[-] (book3) -- (author3);
    \node[emptyStyle](textauthor3) at (10,1) {\it Emile Zola};
   \draw[-] (author3) -- (textauthor3);
    \node[emptyStyle](ln3) at (9,2) {\verb?ln = "Paris"?};
   \draw[-] (author3) -- (ln3);
    \node[emptyStyle](title3) at (11,2) {\verb?<title> ... </title>?};
   \draw[-] (book3) -- (title3);
    \node[emptyStyle](texttitle3) at (11,1) {\it L'assomoir};
   \draw[-] (title3) -- (texttitle3);

\end{scope}

%  \begin{pgfonlayer}{background}
%    \filldraw [line width=4mm,join=round,black!20]
%      (parser.north  -| parser.west)  rectangle (serial.south  -| serial.east);
%\end{pgfonlayer}


\end{tikzpicture}
\end{center}


\framebreak

On poursuit notre chemin vers les enfants (si le nom est \texttt{library})\,: 
\texttt{\textcolor{red}{\bf /}\textcolor{DarkViolet}{\bf child::library}/child::book/child::author}

\begin{center}
\begin{tikzpicture}[auto]
\tiny
  \tikzstyle{emptyStyle}=[text width=0.9cm, text centered]
  \tikzstyle{red}=[rounded corners=1mm,thick,draw=red!75,fill=red!20,text width=0.9cm, text centered]
  \tikzstyle{DarkViolet}=[rounded corners=1mm,thick,draw=DarkViolet!75,fill=DarkViolet!20,minimum size=7mm,text width=0.8cm, text centered]
 \tikzstyle{every label}=[black]

  \begin{scope}
    \node[red](r) at (6,5) {Root};
   \node[emptyStyle](comment) at (1,4) {\texttt{\verb?<!-- ceci ... >?}};
   \draw[-] (r) -- (comment);
    %\node[emptyStyle](entete) at (3,4) {\texttt{$<?xml$ version $...>$}};
   %\draw[-] (r) -- (entete);
    \node[DarkViolet](library) at (6,4) {\verb?<library> ... </library>?};
   \draw[-] (r) -- (library);

    \node[emptyStyle](book1) at (3,3) {\verb?<book> ... </book>?};
   \draw[-] (library) -- (book1);
    \node[emptyStyle](book2) at (6,3) {\verb?<book> ... </book>?};
   \draw[-] (library) -- (book2);
    \node[emptyStyle](book3) at (10,3) {\verb?<book> ... </book>?};
   \draw[-] (library) -- (book3);

    \node[emptyStyle](year1) at (2,3) {\verb?year = "1857"?};
   \draw[-] (book1) -- (year1);
    \node[emptyStyle](author1) at (2,2) {\verb?<author> ... </author>?};
   \draw[-] (book1) -- (author1);
    \node[emptyStyle](textauthor1) at (2,1) {\it Charles Baudelaire};
   \draw[-] (author1) -- (textauthor1);
    \node[emptyStyle](title1) at (3,2) {\verb?<title> ... </title>?};
   \draw[-] (book1) -- (title1);
    \node[emptyStyle](texttitle1) at (3,1) {\it Les fleurs du mal};
   \draw[-] (title1) -- (texttitle1);

    \node[emptyStyle](year2) at (5,3) {\verb?year = "1862"?};
   \draw[-] (book2) -- (year2);
    \node[emptyStyle](author2) at (6,2) {\verb?<author> ... </author>?};
   \draw[-] (book2) -- (author2);
    \node[emptyStyle](textauthor2) at (6,1) {\it Victor Hugo};
   \draw[-] (author2) -- (textauthor2);
    \node[emptyStyle](ln2) at (5,2) {\verb?ln = "Besancon"?};
   \draw[-] (author2) -- (ln2);
    \node[emptyStyle](title2) at (7,2) {\verb?<title> ... </title>?};
   \draw[-] (book2) -- (title2);
    \node[emptyStyle](texttitle2) at (7,1) {\it Les miserables};
   \draw[-] (title2) -- (texttitle2);

    \node[emptyStyle](year3) at (9,3) {\verb?year = "1877"?};
   \draw[-] (book3) -- (year3);
    \node[emptyStyle](author3) at (10,2) {\verb?<author> ... </author>?};
   \draw[-] (book3) -- (author3);
    \node[emptyStyle](textauthor3) at (10,1) {\it Emile Zola};
   \draw[-] (author3) -- (textauthor3);
    \node[emptyStyle](ln3) at (9,2) {\verb?ln = "Paris"?};
   \draw[-] (author3) -- (ln3);
    \node[emptyStyle](title3) at (11,2) {\verb?<title> ... </title>?};
   \draw[-] (book3) -- (title3);
    \node[emptyStyle](texttitle3) at (11,1) {\it L'assomoir};
   \draw[-] (title3) -- (texttitle3);

\end{scope}

%  \begin{pgfonlayer}{background}
%    \filldraw [line width=4mm,join=round,black!20]
%      (parser.north  -| parser.west)  rectangle (serial.south  -| serial.east);
%\end{pgfonlayer}


\end{tikzpicture}
\end{center}

\framebreak

Continuons à descendre, pas à pas.
\smallskip 
\texttt{\textcolor{red}{\bf /}\textcolor{DarkViolet}{\bf child::library}\textcolor{DarkOliveGreen}{\bf /child::book}/child::author}

\begin{center}
\begin{tikzpicture}[auto]
\tiny
  \tikzstyle{emptyStyle}=[text width=0.9cm, text centered]
  \tikzstyle{red}=[rounded corners=1mm,thick,draw=red!75,fill=red!20,text width=0.9cm, text centered]
  \tikzstyle{DarkViolet}=[rounded corners=1mm,thick,draw=DarkViolet!75,fill=DarkViolet!20,minimum size=7mm,text width=0.8cm, text centered]
  \tikzstyle{green}=[rounded corners=1mm,thick,draw=DarkOliveGreen!75,fill=DarkOliveGreen!20,minimum size=7mm,text width=0.8cm, text centered]
 \tikzstyle{every label}=[black]

  \begin{scope}
    \node[red](r) at (6,5) {Root};
   \node[emptyStyle](comment) at (1,4) {\texttt{\verb?<!-- ceci ... >?}};
   \draw[-] (r) -- (comment);
    %\node[emptyStyle](entete) at (3,4) {\texttt{$<?xml$ version $...>$}};
   %\draw[-] (r) -- (entete);
    \node[DarkViolet](library) at (6,4) {\verb?<library> ... </library>?};
   \draw[-] (r) -- (library);

    \node[green](book1) at (3,3) {\verb?<book> ... </book>?};
   \draw[-] (library) -- (book1);
    \node[green](book2) at (6,3) {\verb?<book> ... </book>?};
   \draw[-] (library) -- (book2);
    \node[green](book3) at (10,3) {\verb?<book> ... </book>?};
   \draw[-] (library) -- (book3);

    \node[emptyStyle](year1) at (2,3) {\verb?year = "1857"?};
   \draw[-] (book1) -- (year1);
    \node[emptyStyle](author1) at (2,2) {\verb?<author> ... </author>?};
   \draw[-] (book1) -- (author1);
    \node[emptyStyle](textauthor1) at (2,1) {\it Charles Baudelaire};
   \draw[-] (author1) -- (textauthor1);
    \node[emptyStyle](title1) at (3,2) {\verb?<title> ... </title>?};
   \draw[-] (book1) -- (title1);
    \node[emptyStyle](texttitle1) at (3,1) {\it Les fleurs du mal};
   \draw[-] (title1) -- (texttitle1);

    \node[emptyStyle](year2) at (5,3) {\verb?year = "1862"?};
   \draw[-] (book2) -- (year2);
    \node[emptyStyle](author2) at (6,2) {\verb?<author> ... </author>?};
   \draw[-] (book2) -- (author2);
    \node[emptyStyle](textauthor2) at (6,1) {\it Victor Hugo};
   \draw[-] (author2) -- (textauthor2);
    \node[emptyStyle](ln2) at (5,2) {\verb?ln = "Besancon"?};
   \draw[-] (author2) -- (ln2);
    \node[emptyStyle](title2) at (7,2) {\verb?<title> ... </title>?};
   \draw[-] (book2) -- (title2);
    \node[emptyStyle](texttitle2) at (7,1) {\it Les miserables};
   \draw[-] (title2) -- (texttitle2);

    \node[emptyStyle](year3) at (9,3) {\verb?year = "1877"?};
   \draw[-] (book3) -- (year3);
    \node[emptyStyle](author3) at (10,2) {\verb?<author> ... </author>?};
   \draw[-] (book3) -- (author3);
    \node[emptyStyle](textauthor3) at (10,1) {\it Emile Zola};
   \draw[-] (author3) -- (textauthor3);
    \node[emptyStyle](ln3) at (9,2) {\verb?ln = "Paris"?};
   \draw[-] (author3) -- (ln3);
    \node[emptyStyle](title3) at (11,2) {\verb?<title> ... </title>?};
   \draw[-] (book3) -- (title3);
    \node[emptyStyle](texttitle3) at (11,1) {\it L'assomoir};
   \draw[-] (title3) -- (texttitle3);

\end{scope}

%  \begin{pgfonlayer}{background}
%    \filldraw [line width=4mm,join=round,black!20]
%      (parser.north  -| parser.west)  rectangle (serial.south  -| serial.east);
%\end{pgfonlayer}


\end{tikzpicture}
\end{center}


\framebreak

On arrive au bout\,!
\smallskip 

\texttt{\textcolor{red}{\bf /}\textcolor{DarkViolet}{\bf child::library}\textcolor{DarkOliveGreen}{\bf /child::book}\textcolor{DarkCyan}{\bf  /child::author}}


\begin{center}
\begin{tikzpicture}[auto]
\tiny
  \tikzstyle{emptyStyle}=[text width=0.9cm, text centered]
  \tikzstyle{red}=[rounded corners=1mm,thick,draw=red!75,fill=red!20,text width=0.9cm, text centered]
  \tikzstyle{DarkViolet}=[rounded corners=1mm,thick,draw=DarkViolet!75,fill=DarkViolet!20,minimum size=7mm,text width=0.8cm, text centered]
  \tikzstyle{green}=[rounded corners=1mm,thick,draw=DarkOliveGreen!75,fill=DarkOliveGreen!20,minimum size=7mm,text width=0.8cm, text centered]
  \tikzstyle{cyan}=[rounded corners=1mm,thick,draw=DarkCyan!75,fill=DarkCyan!20,minimum size=7mm,text width=0.8cm, text centered]
 \tikzstyle{every label}=[black]

  \begin{scope}
    \node[red](r) at (6,5) {Root};
   \node[emptyStyle](comment) at (1,4) {\texttt{\verb?<!-- ceci ... >?}};
   \draw[-] (r) -- (comment);
    %\node[emptyStyle](entete) at (3,4) {\texttt{$<?xml$ version $...>$}};
   %\draw[-] (r) -- (entete);
    \node[DarkViolet](library) at (6,4) {\verb?<library> ... </library>?};
   \draw[-] (r) -- (library);

    \node[green](book1) at (3,3) {\verb?<book> ... </book>?};
   \draw[-] (library) -- (book1);
    \node[green](book2) at (6,3) {\verb?<book> ... </book>?};
   \draw[-] (library) -- (book2);
    \node[green](book3) at (10,3) {\verb?<book> ... </book>?};
   \draw[-] (library) -- (book3);

    \node[emptyStyle](year1) at (2,3) {\verb?year = "1857"?};
   \draw[-] (book1) -- (year1);
    \node[cyan](author1) at (2,2) {\verb?<author> ... </author>?};
   \draw[-] (book1) -- (author1);
    \node[emptyStyle](textauthor1) at (2,1) {\it Charles Baudelaire};
   \draw[-] (author1) -- (textauthor1);
    \node[emptyStyle](title1) at (3,2) {\verb?<title> ... </title>?};
   \draw[-] (book1) -- (title1);
    \node[emptyStyle](texttitle1) at (3,1) {\it Les fleurs du mal};
   \draw[-] (title1) -- (texttitle1);

    \node[emptyStyle](year2) at (5,3) {\verb?year = "1862"?};
   \draw[-] (book2) -- (year2);
    \node[cyan](author2) at (6,2) {\verb?<author> ... </author>?};
   \draw[-] (book2) -- (author2);
    \node[emptyStyle](textauthor2) at (6,1) {\it Victor Hugo};
   \draw[-] (author2) -- (textauthor2);
    \node[emptyStyle](ln2) at (5,2) {\verb?ln = "Besancon"?};
   \draw[-] (author2) -- (ln2);
    \node[emptyStyle](title2) at (7,2) {\verb?<title> ... </title>?};
   \draw[-] (book2) -- (title2);
    \node[emptyStyle](texttitle2) at (7,1) {\it Les miserables};
   \draw[-] (title2) -- (texttitle2);

    \node[emptyStyle](year3) at (9,3) {\verb?year = "1877"?};
   \draw[-] (book3) -- (year3);
    \node[cyan](author3) at (10,2) {\verb?<author> ... </author>?};
   \draw[-] (book3) -- (author3);
    \node[emptyStyle](textauthor3) at (10,1) {\it Emile Zola};
   \draw[-] (author3) -- (textauthor3);
    \node[emptyStyle](ln3) at (9,2) {\verb?ln = "Paris"?};
   \draw[-] (author3) -- (ln3);
    \node[emptyStyle](title3) at (11,2) {\verb?<title> ... </title>?};
   \draw[-] (book3) -- (title3);
    \node[emptyStyle](texttitle3) at (11,1) {\it L'assomoir};
   \draw[-] (title3) -- (texttitle3);

\end{scope}

%  \begin{pgfonlayer}{background}
%    \filldraw [line width=4mm,join=round,black!20]
%      (parser.north  -| parser.west)  rectangle (serial.south  -| serial.east);
%\end{pgfonlayer}


\end{tikzpicture}
\end{center}

\end{frame}



\begin{frame}{L'arbre XML vu par XPath}

\begin{center}
\begin{tikzpicture}[auto]
\tiny
  \tikzstyle{emptyStyle}=[text width=0.9cm, text centered]
  \tikzstyle{red}=[rounded corners=1mm,thick,draw=red!75,fill=red!20,text width=0.9cm, text centered]
  \tikzstyle{DarkViolet}=[rounded corners=1mm,thick,draw=DarkViolet!75,fill=DarkViolet!20,minimum size=7mm,text width=0.8cm, text centered]
  \tikzstyle{green}=[rounded corners=1mm,thick,draw=DarkOliveGreen!75,fill=DarkOliveGreen!20,minimum size=7mm,text width=0.8cm, text centered]
  \tikzstyle{cyan}=[rounded corners=1mm,thick,draw=DarkCyan!75,fill=DarkCyan!20,minimum size=7mm,text width=0.8cm, text centered]
 \tikzstyle{every label}=[black]
  \tikzstyle{HotPink}=[rounded corners=1mm,thick,draw=HotPink!75,fill=HotPink!20,minimum size=7mm,text width=0.8cm, text centered]
  \tikzstyle{Olive}=[rounded corners=1mm,thick,draw=Olive!75,fill=Olive!20,minimum size=7mm,text width=0.9cm, text centered]
  \tikzstyle{indiaRed}=[rounded corners=1mm,thick,draw=indiaRed!75,fill=indiaRed!20,minimum size=7mm,text width=0.9cm, text centered]

  \begin{scope}
    \node[red](r) at (6,5) {Root};  
    \node[emptyStyle,text width=1cm](rootnode) at (9,5) {\textcolor{red}{\textit{Root node} (pas un \textit{element node})}};  
   \draw[->,dashed,draw=red] (rootnode) -- (r);
  % \node[indiaRed, label={\textcolor{indiaRed}{\textit{Comment node}}}](comment) at (1,4) {\texttt{\verb?<!-- ceci ... >?}};
   \draw[-] (r) -- (comment);
    %\node[emptyStyle](entete) at (3,4) {\texttt{$<?xml$ version $...>$}};
   %\draw[-] (r) -- (entete);
    \node[emptyStyle,text width=1cm](elementnode) at (10,4) {\textcolor{HotPink}{\textit{Element nodes}}};  
    \node[HotPink](library) at (6,4) {\verb?<library> ... </library>?};
   \draw[-] (r) -- (library);
   \draw[->,dashed,draw=HotPink] (elementnode) -- (library);

    \node[HotPink](book1) at (3,3) {\verb?<book> ... </book>?};
   \draw[-] (library) -- (book1);
   \draw[->,dashed,draw=HotPink] (elementnode) -- (book1);
    \node[HotPink](book2) at (6,3) {\verb?<book> ... </book>?};
   \draw[-] (library) -- (book2);
   \draw[->,dashed,draw=HotPink] (elementnode) -- (book2);
    \node[HotPink](book3) at (10,3) {\verb?<book> ... </book>?};
   \draw[-] (library) -- (book3);

    \node[cyan](year1) at (2,3) {\verb?year = "1857"?};
   \draw[-] (book1) -- (year1);
    \node[HotPink](author1) at (2,2) {\verb?<author> ... </author>?};
   \draw[-] (book1) -- (author1);
    \node[Olive](textauthor1) at (2,1) {\it Charles Baudelaire};
   \draw[-] (author1) -- (textauthor1);
    \node[HotPink](title1) at (3,2) {\verb?<title> ... </title>?};
   \draw[-] (book1) -- (title1);
    \node[Olive](texttitle1) at (3,1) {\it Les fleurs du mal};
   \draw[-] (title1) -- (texttitle1);

    \node[cyan](year2) at (5,3) {\verb?year = "1862"?};
   \draw[-] (book2) -- (year2);
    \node[HotPink](author2) at (6,2) {\verb?<author> ... </author>?};
   \draw[-] (book2) -- (author2);
    \node[Olive](textauthor2) at (6,1) {\it Victor Hugo};
   \draw[-] (author2) -- (textauthor2);
    \node[cyan](ln2) at (5,2) {\verb?ln = "Besancon"?};
   \draw[-] (author2) -- (ln2);
    \node[HotPink](title2) at (7,2) {\verb?<title> ... </title>?};
   \draw[-] (book2) -- (title2);
    \node[Olive](texttitle2) at (7,1) {\it Les miserables};
   \draw[-] (title2) -- (texttitle2);


    \node[cyan](year3) at (9,3) {\verb?year = "1877"?};
   \draw[-] (book3) -- (year3);
    \node[HotPink](author3) at (10,2) {\verb?<author> ... </author>?};
   \draw[-] (book3) -- (author3);
    \node[Olive](textauthor3) at (10,1) {\it Emile Zola};
   \draw[-] (author3) -- (textauthor3);
    \node[cyan](ln3) at (9,2) {\verb?ln = "Paris"?};
   \draw[-] (author3) -- (ln3);
    \node[HotPink](title3) at (11,2) {\verb?<title> ... </title>?};
   \draw[-] (book3) -- (title3);
    \node[Olive](texttitle3) at (11,1) {\it L'assomoir};
   \draw[-] (title3) -- (texttitle3);
    \node[emptyStyle,text width=1cm](textnode) at (4,0) {\textcolor{Olive}{\textit{Text nodes}}};  
   \draw[->,dashed,draw=Olive] (textnode) -- (textauthor1);
   \draw[->,dashed,draw=Olive] (textnode) -- (texttitle1);
   \draw[->,dashed,draw=Olive] (textnode) -- (textauthor2);
    \node[emptyStyle,text width=1cm](attnode) at (8,0) {\textcolor{cyan}{\textit{Attribute nodes}}};  
   \draw[->,dashed,draw=cyan] (attnode) -- (ln2);
   \draw[->,dashed,draw=cyan] (attnode) -- (ln3);
   \draw[->,dashed,draw=cyan] (attnode) -- (year3);
\end{scope}

%  \begin{pgfonlayer}{background}
%    \filldraw [line width=4mm,join=round,black!20]
%      (parser.north  -| parser.west)  rectangle (serial.south  -| serial.east);
%\end{pgfonlayer}

\end{tikzpicture}
\end{center}

%+ \textit{Namespace nodes} + \textit{Processing instruction nodes} \prof{(que nous verrons plus loin dans le cours).}

%NB : XPath ``ne voit pas'' le schéma (doctype). \prof{cad qu'on ne peut pas l'interroger}

\end{frame}




\begin{frame}{Qu'est-ce qu'une \alert{étape} XPath}

Une \alert{étape}  est une expression de la forme 
\texttt{\textcolor{indiaRed}{axis}::\textcolor{DarkOliveGreen}{node-test}\textcolor{cyan}{[predicate]}
}\\


\begin{description}
\item[\texttt{\textcolor{indiaRed}{axis}}] est la direction du pas (\texttt{\textcolor{indiaRed}{child}}  par défaut)\\
Ex : \texttt{\textcolor{indiaRed}{child}} ``descend'' sur les enfants, \texttt{\textcolor{indiaRed}{attribut}} ``descend'' sur les attributs.
\item[\texttt{\textcolor{DarkOliveGreen}{node-test}}] est le type de nœud à retenir\\
Ex : \texttt{\textcolor{DarkOliveGreen}{book}} teste si l'élément atteint (p.e. nœud ou attribut) est un élément nommé \texttt{book}.
Si \texttt{\textcolor{DarkOliveGreen}{*}},  tout nom.
\item[\texttt{\textcolor{cyan}{[predicate]}}] (optionnel) filtre les nœuds selon une expression booléenne \\
	Ex : \texttt{\textcolor{cyan}{[attribute::year='1857']}} teste si le nœud a un attribut \texttt{year} dont la valeur est \texttt{1857}.\\
%	Possibilité d'utiliser des fonctions prédéfinies dans XPath dans les prédicats, par exemple : \texttt{\textcolor{cyan}{[position()=2]}} teste si le nœud est en seconde position, \texttt{\textcolor{cyan}{[position()=last()]}} teste si le nœud est en derniére position.\\
\end{description}

\medskip

Dans \texttt{child::library}, \texttt{child} est l'axe, \texttt{library} le test,  pas de prédicat. 

\begin{remark}
L'axe \texttt{child} est l'axe par défaut, donc on peut simplement écrire \texttt{library}.
\end{remark}
\end{frame}


\begin{frame}{Expression de chemin XPath (définition simplifiée)}

Une expression de chemin est simplement une séquence d'étapes\,:

\vfill

Expression de \alert{chemin absolu} (évaluée à la racine) : 
\texttt{/\textcolor{DarkViolet}{step1}/\textcolor{DarkViolet}{step2}/\textcolor{DarkViolet}{step3}/…/\textcolor{DarkViolet}{stepN}}
\\

\vfill 

Expression de \alert{chemin relative} (évaluée par rapport à un contexte donné) : 
\texttt{\textcolor{DarkViolet}{step1}/\textcolor{DarkViolet}{step2}/\textcolor{DarkViolet}{step3}/…/\textcolor{DarkViolet}{stepN}}
\\

\end{frame}

\begin{frame}{Quelques exemples simples: à vous de jouer...}

\vspace{-0.3cm}

\begin{footnotesize}

\lstinputlisting[basicstyle=\tiny, language=XML,breaklines=true, caption={Fichier limonade.xml}]{codes/limonade.xml}

\vspace{-0.3cm}

\begin{exoblock}

évaluer (sur papier) les requêtes suivantes, sur le fichier limonade.xml se trouvant ci-dessus.

\noindent \texttt{1) /child::inventory \\
2) /child::inventory/child::drink \\
3) /child::inventory/child::drink/child::*\\
4) /child::inventory/child::drink/child::lemonade\\
5) /child::inventory/child::drink/child::*/child::price\\
6) /child::inventory/child::drink/child::price\\
}

Vérifiez avec eXist ou BaseX en chargeant le fichier \texttt{limonade.xml}.
\end{exoblock}
\end{footnotesize}

\end{frame}


\begin{frame}{Interprétation d'un chemin de localisation}

\begin{block}{Intuition}
Un chemin de localisation XPath est une suite d'étapes (steps), chaque étape comprenant éventuellement des tests 
sur les n{\oe}uds atteints pour en éliminer certains.
\end{block}

\medskip
A chaque étape\,:
\begin{enumerate}
  \item on consisère un \alert{n{\oe}ud courant}, pris comme point de départ\,;
   \item on ``sélectionne'' à partir du n{\oe}ud courant un ensemblre d'autres n{\oe}uds\,;
   \item on applique à chacun les tests de n{\oe}ud\,;
   \item on applique les prédicats (optionnels)\,;
   \item enfin on prend chacun des n{\oe}uds restant comme n{\oe}ud courant pour l'étape suivante.
\end{enumerate}
\end{frame}


\begin{frame}{Contexte d'évaluation}

Une expression de chemin est évaluée par rapport à un \textit{contexte}.

\bigskip

Un contexte $[< N1,N2,...,Nn >,Nc]$ consiste en \cite{WDMD} :
\begin{itemize}
\item Une liste de n{\oe}uds $< N1,N2,...,Nn >$ de l'arbre XML
\item Un n{\oe}ud courant $Nc$ appartenant à la liste de n{\oe}uds
\end{itemize}

\medskip

\begin{block}{Concernant le contexte}
La taille du contexte $n$ indique la taille de la liste de n{\oe}uds. On peut la récupérer à l'aide de la fonction \verb?last()?;

\medskip

La position du n{\oe}ud contexte $c \in [1, n]$ indique la position du n{\oe}ud courant au sein de la liste de n{\oe}ud du contexte.
On peut la récupérer à l'aide de la fonction \verb?position()?.
\end{block}


\end{frame}


\begin{frame}[allowframebreaks]{La suite des contextes, jusqu'au résultat}

\begin{center}
\figSized{images/evalxpath.jpg}{évaluation d'une requ\^ete XPath}{0.6}
\end{center}

Le résultat final est l'union des contextes finaux. 

\framebreak
Un pas $step_i$ est évalué par rapport au contexte produit par l'évaluation de $step_{i-1}$. 

\begin{description}
 \item[Pour $i=1$ (premier pas)] si l'expression est \alert{absolue} alors le contexte initial est $[<r>,r]$
       où $r$ est le n{\oe}ud racine. Si l'expression est \alert{relative} alors le contexte est défini par l'environnement.
\item[Pour $i > 1$] si $\mathcal{N} =< N1,N2,...,Nn >$ est le résultat de l'évaluation du pas $step_{i-1}$ 
     alors $step_i$  est successivement évalué par rapport au contexte $[\mathcal{N},Nj]$, pour chaque $j \in [1,n]$.
\end{description}
Le résultat de l'évaluation est l'\textcolor{orange}{\textbf{ensemble}} de n{\oe}uds obtenu après l'évaluation du dernier pas.

\medskip

%à un niveau de détail plus fin, pour chaque pas de la forme \texttt{\textcolor{indiaRed}{axis}::\textcolor{DarkOliveGreen}{node-test}\textcolor{cyan}{[predicate]}}, le contexte passe (est transformé) successivement par l'axe (\texttt{\textcolor{indiaRed}{axis}}), puis le test de n{\oe}ud (\texttt{\textcolor{DarkOliveGreen}{node-test}}), puis le prédicat (\texttt{\textcolor{cyan}{[predicate]}}). 

\end{frame}

